\documentclass[12pt,a4paper]{article}
\usepackage[UTF8]{ctex}
\usepackage{geometry}
\usepackage{amsmath,amssymb}
\usepackage{enumitem}
\usepackage{fancyhdr}
\usepackage{setspace}
\usepackage{hyperref}

\geometry{left=2.5cm,right=2.5cm,top=2.5cm,bottom=2.5cm}
\setlength{\headheight}{15pt}
\setstretch{1.5}
\raggedbottom
\setlength{\emergencystretch}{3em}
\hfuzz=100pt
\hbadness=10000
\hypersetup{hidelinks}

\pagestyle{fancy}
\fancyhf{}
\fancyhead[C]{说明书}
\fancyfoot[C]{\thepage}

\begin{document}

\begin{center}
    {\Huge\heiti 说明书}
\end{center}

\vspace{1em}

\noindent\textbf{发明名称：}基于硅基光电集成的微型化 QKD 芯片封装与校准技术

\section*{一、技术领域}

本发明涉及硅基光电集成、量子密钥分发芯片、片上波导校准、热应力补偿和光电封装技术领域，尤其涉及一种适用于 DV-QKD 和 CV-QKD 芯片化设备的微型化 QKD 芯片封装结构、校准控制方法、模组、终端设备及存储介质。

更具体地，本发明用于解决手机、手表、便携终端、小型网关等微型终端中硅光 QKD 芯片所面临的热漂移、应力漂移、振动扰动和校准污染问题。

\section*{二、背景技术}

现有量子密钥分发设备多采用机架式或台架式结构，能够通过外部温控和光学仪表完成较精细的标定与校准。然而，当量子通信设备向消费电子、便携终端和模块化安全设备演进时，系统体积、功耗、封装空间和环境扰动显著受限，传统台架式校准方式难以继续适用。

在芯片级 QKD 场景下，主光路中的相位、偏振、消光比和耦合状态受温度漂移、装配残余应力、封装材料蠕变、跌落振动以及终端工作热源的共同影响。若直接连续注入高功率参考光用于校准，则会污染量子业务时隙，导致探测背景升高，破坏量子态制备与检测稳定性。

现有技术中，一类方案关注集成 QKD 芯片基本架构，证明芯片化 QKD 的可行性；另一类方案关注通用硅光芯片的稳定化、自校准或可编程光路校准；还有一类方案关注封装层面的热管理与应力控制。但是，这些方案通常缺少面向 QKD 低功率、低噪声要求的封装-校准一体化设计，难以直接适配手机、可穿戴设备等快速温漂和强机械扰动环境。

因此，亟需一种面向终端设备的微型化 QKD 芯片封装与校准技术，以使芯片级 QKD 能够在终端环境中稳定工作。

\section*{三、发明内容}

\subsection*{（一）发明目的}

本发明的目的在于提供一种基于硅基光电集成的微型化 QKD 芯片封装与校准技术，以解决现有技术中芯片级 QKD 难以在手机、手表等微型终端环境下稳定运行的问题。

\subsection*{（二）技术方案}

为实现上述目的，本发明提供如下技术方案。

一种面向量子密钥分发芯片的硅基光电集成封装结构，包括：
\begin{enumerate}[leftmargin=2em]
    \item 硅光子芯片，所述硅光子芯片上形成承载量子业务的主光路、与主光路共模布置的参考波导、用于表征温度漂移的片上传感单元、用于表征封装应力或应变的传感单元以及用于补偿相位、偏置或损耗的执行单元；
    \item 封装基座，所述硅光子芯片固定于低热膨胀基座或低应力基座上；
    \item 应力释放结构，设置于芯片周围或封装互连区域，用于减小外壳形变、装配残余应力和温度循环对核心光路区的影响；
    \item 控制电路，所述控制电路基于参考波导和传感单元的观测结果生成补偿控制量；
    \item 光学输入输出耦合单元，用于将量子业务光和参考光耦合至芯片或从芯片输出。
\end{enumerate}

在一个实施方式中，所述参考波导与主光路采用并排、镜像或环绕布置，并与热源器件保持预设热耦合和热隔离关系，以提高温漂与应力漂移的可观测性。

在一个实施方式中，所述传感单元包括功率抽头与监测光电二极管、参考马赫-曾德干涉仪、参考微环谐振器、温度传感器、应变计和应力释放梁读出结构中的至少一种。

在一个实施方式中，所述执行单元包括加热器阵列、相位调节器、IQ 偏置调节单元、可调衰减器和光开关中的至少一种。

一种基于上述封装结构的 QKD 芯片校准控制方法，包括：
\begin{enumerate}[leftmargin=2em]
    \item 将量子业务组织为包含量子业务时隙和静默校准时隙的帧结构或超级帧结构；
    \item 在静默校准时隙中以低占空比、低功率向参考波导和选定主光路段注入校准信号；
    \item 采集参考波导、温度传感单元和应力传感单元的观测结果；
    \item 根据所述观测结果估计温度漂移分量和应力漂移分量，并生成对执行单元的补偿控制量；
    \item 在后续量子业务时隙中使主光路按照补偿后的工作点运行，以维持量子态制备或检测稳定性。
\end{enumerate}

在一个实施方式中，系统建立状态更新关系：
\[
\mathbf{x}_{k+1}=\mathbf{F}\mathbf{x}_k+\mathbf{G}\mathbf{u}_k+\mathbf{w}_k,
\]
\[
\mathbf{y}_k=\mathbf{H}\mathbf{x}_k+\mathbf{v}_k,
\]
其中，$\mathbf{x}_k$ 表示温度、应力、相位和偏置相关状态量，$\mathbf{u}_k$ 表示执行单元控制量，$\mathbf{y}_k$ 表示传感器观测量。系统可采用卡尔曼滤波、查找表补偿、模型预测控制、在线辨识或数据驱动回归等方式完成状态估计和补偿。

在一个实施方式中，参考光路与量子业务光路之间通过时分隔离、波分隔离、片上滤波、封装内滤波和探测窗口门控中的至少一种实现隔离，以降低参考光对量子业务的背景污染。

本发明还提供一种微型化 QKD 模组、一种终端设备、一种电子设备以及一种计算机可读存储介质，用于执行上述结构和方法。

\subsection*{（三）有益效果}

与现有技术相比，本发明至少具有如下有益效果：
\begin{enumerate}[leftmargin=2em]
    \item 通过参考波导与主光路共模布置，提高对温漂和应力漂移的可观测性；
    \item 通过静默时隙校准机制，实现校准和量子业务共存，降低参考信号污染；
    \item 通过封装结构、片上传感和控制电路协同设计，缩短反馈链路并降低漂移源；
    \item 通过温度和应力联合估计，提高在快速热扰动和机械扰动下的稳定性；
    \item 适用于手机、智能手表、便携网关等多类终端场景，具有较高工程落地价值。
\end{enumerate}

\section*{四、附图说明}

为了更清楚地说明本发明实施例或现有技术中的技术方案，下面对实施例中所需要使用的附图作简单说明。所述附图仅为示意性附图。

\begin{enumerate}[leftmargin=2em]
    \item 图1为本发明微型化 QKD 芯片封装结构总体示意图；
    \item 图2为本发明主光路、参考波导、片上传感单元与执行单元的共模布局示意图；
    \item 图3为本发明静默校准时隙与量子业务时隙交织的帧结构示意图；
    \item 图4为本发明片上传感、状态估计和执行器反馈控制闭环示意图；
    \item 图5为本发明在手机、可穿戴设备和小型网关中的实施例示意图。
\end{enumerate}

\section*{五、具体实施方式}

下面结合附图和具体实施方式对本发明作进一步详细说明。应当理解，这里描述的实施方式仅用于解释本发明，并不用于限定本发明。

\subsection*{（一）术语定义}

在本说明书中：
\begin{itemize}[leftmargin=2em]
    \item ``主光路''表示承载量子业务信号的片上光路；
    \item ``参考波导''表示与主光路在热环境和几何布局上至少部分共模、用于观测漂移状态的辅助光路；
    \item ``静默校准时隙''表示在不承载量子业务的情况下执行参考注入和状态观测的短时隙；
    \item ``执行单元''表示用于补偿相位、偏置或损耗的片上或近端器件；
    \item ``控制电路''表示用于采集观测量、执行状态估计和输出补偿控制量的 ASIC、MCU 或其他控制器。
\end{itemize}

\subsection*{（二）差分参考结构实施方式}

如图1和图2所示，在一个实施方式中，硅光子芯片同时布置主光路和参考波导。主光路用于量子态制备或检测，参考波导用于在不显著干扰量子业务的情况下感知热漂移和应力漂移。参考波导与主光路在同一区域并行延伸，从而共享主要共模温漂，但在局部结构上保持可辨识的差分路径，以便推断差分相位偏移。

在一个实施方式中，芯片周围设置对称焊盘、对称粘接区和应力释放槽，以减小封装弯曲或外壳形变对主光路的影响。

\subsection*{（三）片上传感与状态估计实施方式}

如图4所示，在一个实施方式中，系统在主光路关键节点设置功率抽头和监测光电二极管，在参考波导上设置参考干涉结构或参考微环，同时设置温度传感单元和应力传感单元。控制电路读取上述观测量，并按照状态空间模型估计温度漂移分量与应力漂移分量。

在一个实施方式中，控制电路针对毫秒级快速扰动执行小步快速补偿，针对秒级或更长时间尺度的缓慢漂移执行深度重锁定或低频扫描。

\subsection*{（四）静默时隙校准实施方式}

如图3所示，在一个实施方式中，量子业务帧由若干量子业务时隙和若干静默校准时隙组成。在静默校准时隙中，控制电路短时打开参考光路，向参考波导和选定主光路段注入低占空比参考信号，并立即读取观测量。随后，根据估计结果更新执行单元工作点，在下一个量子业务时隙中以新的工作点运行主光路。

在一个实施方式中，当检测到终端进入充电、大功率射频发射、振动或跌落恢复阶段时，控制电路临时提高静默校准时隙密度。

\subsection*{（五）封装协同实施方式}

在一个实施方式中，控制 ASIC 与光子芯片采用近距共封装，以缩短反馈链路延迟。光学输入输出可根据终端形态选择边耦合、光栅耦合或聚合物波导过渡，并在相应区域设置应变隔离结构，以提高长期耦合稳定性。

在一个实施方式中，封装盖板或腔体还可用于减小湿热和颗粒污染对器件漂移的影响。

\subsection*{（六）终端实施例}

如图5所示，在一个手机终端实施例中，QKD 芯片与控制 ASIC 共同封装于微型模组内。由于应用处理器和射频前端的工作，芯片周围温度在几秒内快速变化。系统通过参考微环和应力传感结构检测漂移，在连续两个静默校准时隙内完成状态估计和执行器修正，使下一业务窗口恢复到目标误码率和干涉可见度范围。

在一个可穿戴实施例中，终端空间更紧凑，封装易受弯折与振动影响，因此优选采用更高柔顺度的应力释放结构并提高应力传感权重。

在一个小型网关实施例中，由于供电和散热条件相对稳定，系统可采用更低频率的静默校准，并在业务低谷期执行深度重锁定。

\subsection*{（七）适用范围说明}

本发明适用于 DV-QKD、CV-QKD、MDI-QKD 及其他需要片上稳定量子态制备或检测的芯片化量子通信系统；适用于手机、智能手表、便携终端、小型网关和加密钥匙等多种终端形态。只要采用所述主光路与参考波导协同布局、静默时隙校准、温应力联合估计和封装协同反馈的技术思想，均应落入本发明保护范围。

\end{document}
